\documentclass[letterpaper,11pt]{article}
\usepackage[top=0.8in, bottom=0.8in, left=0.9in, right=0.9in]{geometry}
\usepackage[hidelinks]{hyperref}
\usepackage{lmodern}
\usepackage{parskip}
\usepackage{microtype}
\setlength{\parskip}{6pt}
\setlength{\parindent}{0pt}
\begin{document}
\begin{flushleft}
\textbf{\large Ajay Venkatesh} \\
Brooklyn, NY \\
+1 516 585 9013 \\
\href{mailto:av3855@nyu.edu}{av3855@nyu.edu}
\end{flushleft}
\vspace{10pt}
May 21, 2026
\vspace{6pt}

Hiring Manager \\
San Diego State University Research Foundation

\vspace{10pt}

Dear Hiring Manager,

My name is Ajay Venkatesh. I'm finishing my M.S. in Computer Engineering at NYU, with production software experience at PwC in Bengaluru, a summer SDE internship at Amazon, and ongoing on-campus work at NYU's Novel AI Technologies. I'm writing about the Data Scientist (Data Analytics Developer II) role at SDSU Research Foundation.

The role's combination of data analytics development with research-supporting infrastructure is where my engineering work has been concentrated. At NYU's Novel AI Technologies I engineered \textbf{ETL pipelines} transforming \textbf{NoSQL} into \textbf{PostgreSQL} with concurrency control and locking mechanisms, processing hundreds of thousands of records, then powered \textbf{AWS QuickSight} dashboards for real-time analytics. At PwC I owned a distributed backend over \textbf{SQL Server}, writing complex queries and parallelized batch ingestion that materially cut execution time.

On large-scale data analysis: my Big Data graduate project built a \textbf{PySpark / HDFS} pipeline over millions of flight records, training Random Forest and Gradient Boosted Tree models with rigorous EDA, anomaly detection, and feature engineering. My coursework covers \textbf{Machine Learning}, \textbf{Deep Learning}, \textbf{Big Data}, and \textbf{Data Structures}, the foundation for the kind of analytical work the JD describes.

On research-supporting infrastructure: I partnered with researchers at NYU's Novel AI Technologies on an NSF-funded health platform, building reproducible data pipelines and analytics dashboards for ongoing research work. Translating research questions into reliable analytical deliverables is a skill that compounds across grant-funded projects.

Stack-wise: \textbf{SQL}, \textbf{Python}, \textbf{Pandas}, \textbf{NumPy}, \textbf{PySpark}, \textbf{scikit-learn}, plus \textbf{PostgreSQL}, \textbf{MySQL}, \textbf{MongoDB}, and visualization through \textbf{Power BI}, \textbf{Tableau}, and \textbf{QuickSight}. Comfortable across data warehousing, ETL workflows, and statistical analysis.

What pulls me toward SDSU Research Foundation is the mission. Data analytics development for research grants and university initiatives runs on careful, reproducible work that compounds across studies, and contributing to research infrastructure beyond any single project is meaningful engineering.

I'd welcome the chance to discuss further. Thanks for considering my application.

\vspace{12pt}
Sincerely, \\
Ajay Venkatesh

\end{document}